\documentclass[11pt]{article}
\usepackage[margin=24mm]{geometry}
\usepackage{amsmath,amssymb,amsthm}
\setlength{\emergencystretch}{2em}
\newtheorem{theorem}{Theorem}\newtheorem{lemma}[theorem]{Lemma}
\newcommand{\F}{\mathbb F}\newcommand{\Sp}{\operatorname{Sp}}\newcommand{\diag}{\operatorname{diag}}\newcommand{\Tr}{\operatorname{Tr}}\newcommand{\Om}{\Omega}
\title{Finite kernels of full-rank orthogonal inverters in characteristic two}
\author{Human Compression project}\date{}
\begin{document}\maketitle
Let $q$ be even, let $m\ge2$, let $V$ be a nonsingular quadratic space
of dimension $2m$, and let $t\in O(V)$ have $(t-I)^2=0$ and rank $m$.
Write $B$ for the polar form and assume $B(v,tv)$ is not identically zero.
These are the full-rank types $c_m$ for even $m$, and $b_m$ for odd $m$
in Aschbacher--Seitz Sections7--8. The result below also covers $t\notin\Om(V)$:
the centralizer is always taken in $\Om(V)$. Put $\nu_2(H)=|\{h:h^2=1\}|$.

\begin{theorem}\label{main}
There are surjective homomorphisms with the indicated finite kernel orders:
\[
\begin{array}{c|c|c}
 &\text{image of }C_{\Om(V)}(t)&\text{kernel order}\\\hline
m\text{ odd}&\Sp_{m-1}(q)&q^{(m^2-m)/2}\\
m\text{ even}&\Sp_{m-2}(q)&q^{(m^2+m-2)/2}.
\end{array}
\]
Here $\Sp_0$ is trivial. Consequently the centralizer's $\nu_2$ is at most
the kernel order times $\nu_2$ of the displayed symplectic group.
In particular, when $m=4$,
\[
                  \nu_2(C_{\Om(V)}(t))\le q^9q^2=q^{11}.
\]
No order is inferred from an algebraic dimension.
\end{theorem}

\section{A direct finite fibre calculation}
The nonalternating full-rank symplectic normal form from AS7.6--7.7
permits coordinates $x,y\in\F_q^m$ in which
\[
 t(x,y)=(x+y,y),\qquad
 B((x,y),(x',y'))=x^ty'+y^tx'.
\]
Equivalently, the nonsingular nonalternating symmetric matrix defining
the square-zero part is congruent to $I_m$ over the perfect field.
Let $s=(1,\ldots,1)^t$. A quadratic form with this polar form differs
from $x^ty$ by squares of linear forms. Its invariance under $t$ forces
its $x$-part to be $(s^tx)^2$. Thus, for some column $b$, it is exactly
\begin{equation}\label{Q}
 Q(x,y)=(s^tx)^2+x^ty+b^ty^{[2]},
\end{equation}
where $[2]$ denotes entrywise squaring.

Block commutation and the symplectic equation give every symplectic
centralizer element uniquely as
\[
 g_{A,D}(x,y)=(Ax+DAy,Ay),\qquad AA^t=I_m,\quad D=D^t.
\]
Put $J_m=\{A:AA^t=I_m\}$. These matrices fix $s$, since
$v^tv=(s^tv)^2$ and squaring is injective. Substituting $w=Ay$ in
\eqref{Q} proves that $g_{A,D}$ is orthogonal precisely when
\begin{equation}\label{F}
 F(D):=\diag(D)+(Ds)^{[2]}=b+A^{[2]}b.
\end{equation}
This is an equation of additive groups, not an asserted $\F_q$-linear map.

We count its fibres. Put $d=m(m-1)/2$. If $c_i$ is the sum of the
off-diagonal entries in row $i$, then $\sum c_i=0$. The map from the
$d$ off-diagonal entries to $(c_i)$ is onto that hyperplane: entries in
the last column prescribe its first $m-1$ coordinates. Every fibre has
$q^{d-m+1}$ elements.
Writing $a_i=D_{ii}$, equation $F(D)=0$ becomes
\[
                       a_i^2+a_i=c_i^2.
\]
Over $\F_q$, the map $a\mapsto a^2+a$ has kernel of order two and image
$\ker\Tr_{q/2}$. Since $\Tr(c_i^2)=\Tr(c_i)$, there are
$(q/2)^{m-1}$ possible row-sum vectors, then $q^{d-m+1}$ off-diagonal
choices, and $2^m$ diagonal choices. Therefore
\begin{equation}\label{K}
                         |\ker F|=2q^d.
\end{equation}
Also, writing $a=\sum_iD_{ii}$, we have
$\sum_iF(D)_i=a+a^2$. It follows that the image of $F$ lies in
$\{c:\Tr_{q/2}(\sum c_i)=0\}$. Both sets have size $q^m/2$ by
\eqref{K}, so they are equal. Since $s^tA^{[2]}=s^t$, the right side
of \eqref{F} has coordinate sum zero. Thus every $A\in J_m$ lifts,
and every fibre in $C_{O(V)}(t)\to J_m$ has exactly $2q^d$ elements.

The matrix $g_{I,E_{11}}$ is in this kernel: $F(E_{11})=0$.
It has rank one after subtraction of the identity, so its Dickson
invariant is one. We retain the standard Dickson homomorphism
$O(V)\to C_2$, $g\mapsto\operatorname{rank}(g-I)\bmod2$, whose kernel
is $\Om(V)$. Multiplication by this kernel element toggles that invariant
without changing $A$. Therefore every fibre has exactly $q^d$ elements
in $\Om(V)$. We have proved the surjection
\begin{equation}\label{surjection}
 C_{\Om(V)}(t)\longrightarrow J_m,\qquad |\ker|=q^d.
\end{equation}
This argument pays the full orthogonal component and its index-two
intersection; no connectedness or splitting assumption replaces the count.

\section{The smaller symplectic image}
For odd $m$, $s^ts=1$, and $s^\perp$ is a nonsingular alternating space.
Every member of $J_m$ fixes $s$, and conversely every symplectic isometry
of $s^\perp$ extends by fixing $s$. Hence $J_m=\Sp_{m-1}(q)$.

For even $m$, $s^ts=0$. Choose $e$ with $s^te=1$, whence $e^te=1$,
and let $H=\langle s,e\rangle^\perp$. The form on $H$ is nonsingular
alternating of dimension $m-2$. The action on $s^\perp/\langle s\rangle$
gives a surjection $J_m\to\Sp(H)$, split by fixing $s,e$.
Every element of its kernel, and no other one, has the form
\[
 s\mapsto s,\qquad h\mapsto h+(h^tu)s\ (h\in H),\qquad
 e\mapsto e+u+a s,\qquad u\in H,\ a\in\F_q.
\]
These maps preserve every displayed pairing. Conversely, fixing the
quotient forces the stated shape, and the pairing with $e$ determines
its linear functional. Thus the kernel has exactly $q^{m-1}$ elements.
Combining with \eqref{surjection} proves the exact kernel orders in
Theorem~\ref{main}.
These kernels have orders which are powers of two. They are therefore
unipotent matrix groups in characteristic two; their orders above are
finite counts, including the orthogonal component factor.

For any surjection $C\to L$ with finite kernel $K$, an element of square
one maps to an element of square one. There are at most $|K|$ such elements
above each image, so $\nu_2(C)\le|K|\nu_2(L)$. In characteristic two,
$\Sp_2(q)=\operatorname{SL}_2(q)$ has exactly $q^2-1$ nonidentity
involutions: each is a rank-one unipotent, with one eigenline, and for
each of the $q+1$ lines there are $q-1$ nontrivial transvections.
Thus $\nu_2(\Sp_2(q))=q^2$, which proves the stated $q^{11}$ bound.

\section{Application to the natural torus inverters}
In the acquired orthogonal torus models, a natural projective involutory
inverter in characteristic two has an isometric lift of square one.
Indeed if a conformal lift has square $aI$, its multiplier is $a$;
rescaling by the unique square root of $a^{-1}$ achieves both conditions.
Projective inversion becomes literal inversion, since the only scalar
isometry in characteristic two is one.

On a negative $2k$-block write $E=\F_{q^{2k}}$, $\bar v=v^{q^k}$ and
$Q(v)=\Tr_{q^k/q}(c_0v\bar v)$, $c_0\ne0$.
The acquired torus degree argument distinguishes this irreducible block
from the others. If its multiplier is $\alpha$, then
$\alpha^{-1}=\bar\alpha$, and an inverting linear map satisfies
$t(\alpha v)=\bar\alpha t(v)$. Since $\F_q[\alpha]=E$, it is exactly
$v\mapsto c\bar v$. Its square is one precisely when $c\bar c=1$.
The small negative plane in the plus model has the same argument,
since its primitive multiplier has order $q+1>2$.
Hilbert90, or the cyclic
multiplicative group directly, conjugates it linearly to $v\mapsto\bar v$,
so its square-zero part has rank $k$. Moreover its polar form satisfies
\[
             B(v,c\bar v)=\Tr_{E/\F_q}(c_0\bar c\,v^2),
\]
which is not identically zero by the nondegenerate trace pairing and
bijectivity of squaring. The minus model therefore has rank $m$ and is
nonalternating. In the plus model the two negative blocks give ranks
$m-1$ and one; restriction to the large block already makes the full
form nonalternating. They therefore have precisely the types $c_m$
($m$ even) and $b_m$ ($m$ odd), and Theorem~\ref{main} applies for either
orthogonal sign. No flag-dimension test is required to eliminate $a_m$.

\paragraph{Source and scope.}
AS1976 Sections7--8 supplies the symplectic normal form and class labels;
the complete relevant passages and the 1978 corrections were read.
The correction to AS8.8(i) concerns the last column of $P$; here the
middle space has dimension zero, and the direct equation \eqref{F}
independently pays all finite fibres. The Dickson convention is retained
from Murray--Roney-Dougal2008, Definition3.1(2), Lemma3.2 and
Definition3.3. The coordinate fibre count
and its use here are new derivations by publication\_recovery; root
requested this interface and proposed its expected exponents before
this proof. These arguments do not pay triality, maximal-overgroup
classification, odd-characteristic fixed groups or the whole L045 claim.

\begin{thebibliography}{9}\small
\bibitem{AS} M.~Aschbacher and G.~M.~Seitz,
\emph{Involutions in Chevalley groups over fields of even order},
Nagoya Math. J. 63 (1976), 1--91, Sections 7--8; corrections,
Nagoya Math. J. 72 (1978), 135--136.
\bibitem{MRD} S.~H.~Murray and C.~M.~Roney-Dougal,
\emph{The spinor norm and homomorphism algorithms for classical groups},
preprint dated 2 May 2008, Section 3.
\end{thebibliography}
\end{document}
